\section{Absolute operator bounds and their infrared cost}
\label{sec:lower-proof}

For the remainder of the paper write
$H_{\Lambda,L}=H_D+H_\gamma+H_I+H_C$, where $H_C$ includes
the four-field Wick Coulomb term, the linear source coupling, and
the nuclear pair energy. The current is
$j_i(x)=:\!\psi^\dagger(x)\mathsf a_i\psi(x)\!:_P$.
These conventions are exactly those of \eqref{eq:model} and
\eqref{pe-field-definition}; no new normalization is introduced.

\begin{lemma}[Local cutoff bounds and noncommutative square completion]
\label{lem:local}
At each point, with $d_L=M_L/V$,
\begin{equation}
 -2d_LI\le\rho(x),j_i(x)\le2d_LI,
 \qquad \sum_{i=1}^3j_i(x)^2\le12d_L^2I.
 \label{eq:local-bounds}
\end{equation}
\end{lemma}
\begin{proof}
The four operators $\psi_a(x)/\sqrt{d_L}$ satisfy normalized
four-mode CAR. Their number operator has spectrum $0,1,2,3,4$,
and the sea has local density $2d_L$. This proves the density bound.
Each $\mathsf a_i$ has two eigenvalues $+1$ and two $-1$, so
diagonalizing it proves the current bound. Its normal-ordering scalar
vanishes by the symmetric momentum cutoff. Squaring the three
separate self-adjoint bounds gives the last inequality.
\end{proof}

\begin{lemma}[Photon square completion]
\label{lem:square}
For a bounded fermion operator $B$ and a commuting photon oscillator
$a$, the following quadratic-form identity holds:
\begin{equation}
 \omega a^\dagger a+a^\dagger B+B^\dagger a
 =\omega(a+\omega^{-1}B)^\dagger(a+\omega^{-1}B)
   -\omega^{-1}B^\dagger B.
 \label{eq:square}
\end{equation}
\end{lemma}
\begin{proof}
Expand the right side. Only commutation between the fermion and
photon tensor factors is used.
\end{proof}
The shifted operators need not be canonical oscillators:
$[a+\omega^{-1}B,a^\dagger+\omega^{-1}B^\dagger]
=I+\omega^{-2}[B,B^\dagger]$. Thus the identity is a valid
lower-bound device, not a simultaneous unitary diagonalization of
noncommuting currents.

Let $H_C^{\rm dens}$ be obtained by replacing the four-field Wick
term with $\alpha D_L(\rho,\rho)/2$. The symmetric Fourier sum
is a positive operator sum. Set $\nu_s=e^{s\Delta}\nu$,
$G_s=e^{s\Delta}G_L$, and $h_s=G_L-G_s$. Direct expansion gives
\begin{align}
 H_C^{\rm dens}
 &=\frac\alpha2D_L(\rho+\nu_s,\rho+\nu_s)
   +\alpha D_L(\rho,\nu-\nu_s)\notag\\
 &\quad+\alpha\sum_{i<j}Z_iZ_jh_{2s}(R_i-R_j)
   -\frac\alpha2G_{2s}(0)\sum_iZ_i^2.
 \label{eq:heat-identity}
\end{align}
This identity uses point nuclei throughout; smoothing is only a
comparison in the proof. Lemma~\ref{lem:heat} and
\eqref{eq:local-bounds} bound the core term below by
$-16\pi\alpha d_LsZ_{\rm tot}$, the nuclear short-range term
below by $-4\pi\alpha sZ_{\rm tot}^2/V$, and the last diagonal
by the stated heat-kernel estimate. The exact neutrality capacity
$Z_{\rm tot}\le2M_L$ therefore gives
\begin{equation}
 H_C^{\rm dens}\ge-24\pi\alpha d_LsZ_{\rm tot}
 -\frac{13\alpha}{4\sqrt2\pi^{3/2}\sqrt s}\sum_iZ_i^2.
 \label{b-density-absolute}
\end{equation}

For completely explicit constants, put $d_*=(1+1/\pi)^3$ and
\begin{align}
 A_\Lambda&=\alpha\Lambda_f
 \left[24\pi d_*Z_{\max}
       +\frac{13}{4\sqrt2\pi^{3/2}}Z_{\max}^2\right],\notag\\
 B_\Lambda&=\frac{364\alpha}{\pi^2}d_*\Lambda_f^4.
 \label{b-absolute-constants}
\end{align}
For $L\ge\Lambda_f^{-1}$, taking $s=\Lambda_f^{-2}$ proves
$H_C^{\rm dens}\ge-A_\Lambda K$.
The exact Wick identity~\eqref{eq:wick} controls the ordering change.
For contractions $A,B$ with norms at most one,
$\|AB-APB-BPA\|\le3$, $\|n_P(X)\|\le8M_L\|X\|$,
and $|\Tr(PA(1-P)B)|\le4M_L$. Hence the norm of the difference
per Coulomb mode is at most $28M_L$. Since
\begin{equation}
 \frac1V\sum_{0<|k|\le2\Lambda_f}\frac{4\pi}{|k|^2}
 =\kappa_{\Lambda,L}\le\frac{26\Lambda_f}{\pi^2},
 \label{eq:coulomb-mode-sum}
\end{equation}
we obtain the full operator estimate
\begin{equation}
 \|H_C-H_C^{\rm dens}\|\le B_\Lambda V,
 \qquad H_C\ge-A_\Lambda K-B_\Lambda V.
 \label{eq:full-coulomb-bound}
\end{equation}
All of these estimates hold over the full admissible Fock sector.

If every coupled photon frequency is at least $k_{\min}>0$, square
completion, transverse polarization projection, and Parseval give
\begin{align}
 H_\gamma+H_I
 &\ge-\frac{e^2}{2V}\sum_{k,\lambda}
       \frac{|\varepsilon_\lambda(k)\cdot\widehat j(k)|^2}{|k|^2}
 \notag\\
 &\ge-\frac{e^2}{2k_{\min}^2}
          \int\sum_i j_i(x)^2\dd x
 \ge-\frac{6e^2d_L^2V}{k_{\min}^2}.
 \label{b-photon-absolute}
\end{align}
Here $|B|^2=B^\dagger B$; polarization and Parseval are operator
quadratic-form identities and require no current commutativity.
For ordinary nonzero torus modes $k_{\min}=2\pi/L$. Thus
\begin{equation}
 H_{\Lambda,L}\ge-A_\Lambda K-B_\Lambda V
  -\frac{6e^2d_*^2\Lambda_f^6}{4\pi^2}VL^2.
 \label{b-lower-L5}
\end{equation}
With a separately imposed fixed infrared cutoff $k_{\min}\ge\kappa>0$
this argument instead gives an extensive lower bound, at the cost of
changing the coupled photon modes. The next section shows that the
$L^5$ loss is real for the unmodified bare regulator with massless photons.
